% =====================================================================
% Matched comparison and ablation results
% =====================================================================

\subsection{Matched tactical baselines}
\label{subsec:comparison}

Only the tactical-guidance module changes across methods. All methods use the
same plant, race-control interface, corridor projection, path and speed
planners, MPCC, and actuator allocator. \emph{Pure Raceline} contains no
opponent-aware tactical decision. \emph{Non-Interactive MPC} predicts the
opponent independently of the ego plan and imposes obstacle constraints on the
result~\cite{kabzan2020amzracing,betz2023autonomous}.
\emph{Rule-based GT} uses a fixed-threshold state machine
($g_{\rm chase}=20$ m, $g_{\rm ot}=8$ m, and $g_{\rm abort}=35$ m).
\emph{iLQGames} computes coupled local responses with iterative
linear-quadratic game approximations~\cite{fridovich2020ilqgames}. Outputs from
all baselines are projected onto the same side, speed-cap, and corridor
interface so that the reported differences concern tactical reasoning rather
than low-level execution.

\begin{table}[!tbp]
\centering
\footnotesize
\setlength{\tabcolsep}{2.3pt}
\renewcommand{\arraystretch}{1.08}
\caption{Matched comparison over 1000 closed-loop HIL laps per method. Arrows
indicate the preferred direction.}
\label{tab:comparison}
\begin{tabular}{@{}lccccc@{}}
\toprule
Metric & Prop. & PR & NI-MPC & RB-GT & iLQG \\
\midrule
OSR [\%] $\uparrow$             & \textbf{94.7} & 38.2 & 28.5 & 51.6 & 63.5 \\
OSR def. [pp] $\downarrow$      & \textbf{0.0} & 56.5 & 66.2 & 43.1 & 31.2 \\
TTO [s] $\downarrow$            & \textbf{4.82} & 7.25 & 11.23 & 9.37 & 7.83 \\
Min. gap [m] $\uparrow$         & \textbf{5.83} & 1.42 & 3.85 & 2.14 & 2.31 \\
Coll. [count] $\downarrow$      & \textbf{0} & 21 & 7 & 48 & 14 \\
Mean TTC [s] $\uparrow$         & \textbf{8.52} & 3.21 & 6.92 & 5.64 & 5.73 \\
Track viol. [\%] $\downarrow$  & \textbf{0.0} & \textbf{0.0} & \textbf{0.0} & \textbf{0.0} & \textbf{0.0} \\
PFR $\uparrow$                  & \textbf{1.000} & 0.988 & 0.998 & 0.987 & 0.995 \\
\bottomrule
\end{tabular}
\vspace{1mm}

\begin{minipage}{\columnwidth}
\scriptsize
\emph{Notes:} Prop., proposed method; PR, Pure Raceline; NI-MPC,
Non-Interactive MPC; RB-GT, Rule-based GT; iLQG, iLQGames. OSR def.
is the percentage-point deficit from Prop. Min. gap is the mean per-lap
minimum inter-vehicle distance. PFR is the fraction of planning cycles with a
feasible SOCP solution. Bold denotes the best observed value; equal track
violation values are jointly bold.
\end{minipage}
\end{table}

Table~\ref{tab:comparison} makes the primary effect explicit through the OSR
and OSR-deficit rows. The proposed tactical corridor policy completed 94.7\%
of the laps, 31.2 percentage points more than the strongest baseline,
iLQGames. It also recorded no collision and retained a feasible path in every
planning cycle.
iLQGames confirms the value of interactive reasoning relative to the
non-interactive optimizer, but its 63.5\% OSR and 14 contacts show that an
online game trajectory alone did not provide the persistence of the learned
corridor interface. Conversely, Non-Interactive MPC maintained the second
largest mean minimum gap and only seven contacts but had the lowest OSR. The
combination is evidence of conservative tactical freezing rather than poor
low-level tracking. The rule-based method produced 48 contacts, showing that
one threshold set did not transfer across the matched curvature and relative-
speed regimes. These contrasts isolate the intended gap: interaction reasoning
must be both opponent-aware and expressed as a persistent, adaptable feasible
region, rather than as an isolated game trajectory, obstacle constraint, or
fixed maneuver threshold.

For transparency, two-sided two-proportion $z$ tests were recomputed from the
machine-readable integer success counts and adjusted across the four baseline comparisons by
the Holm procedure. All count-level differences remained below
$p=10^{-6}$. These tests do not model scenario or training-seed clustering and
are therefore treated as descriptive support; the effect sizes and event
counts in Table~\ref{tab:comparison} carry the primary interpretation.

\subsection{Ablation of the tactical representation}
\label{subsec:ablation}

Four ablations isolate the proposed tactical components while retaining the
same execution layer. \emph{Without game guidance} removes both the IBR action
prior and the auxiliary robust-value target. \emph{Without FSM} bypasses
hysteretic mode and side locking, allowing the policy to reshape the corridor
on every update. \emph{Without learned corridor} replaces the continuous
width, clearance, lateral-bias, and speed-cap outputs by fixed nominal values.
The final variant replaces TQC with Soft Actor-Critic (SAC) while preserving
the network, replay, reward, and other training settings.

\begin{table}[!htbp]
\centering
\footnotesize
\setlength{\tabcolsep}{2.3pt}
\renewcommand{\arraystretch}{1.08}
\caption{Component ablations over 1000 closed-loop HIL laps per variant. Each
variant removes or replaces one tactical component. Arrows indicate the
preferred direction.}
\label{tab:ablation}
\begin{tabular}{@{}lccccc@{}}
\toprule
Metric & Full & No-GG & No-FSM & Fixed-C & SAC \\
\midrule
OSR [\%] $\uparrow$            & \textbf{94.7} & 85.4 & 82.3 & 76.1 & 88.5 \\
OSR loss [pp] $\downarrow$    & \textbf{0.0} & 9.3 & 12.4 & 18.6 & 6.2 \\
TTO [s] $\downarrow$           & \textbf{4.82} & 6.38 & 6.71 & 8.94 & 5.53 \\
Coll. [count] $\downarrow$    & \textbf{0} & 9 & 12 & 31 & 5 \\
Min. gap [m] $\uparrow$       & \textbf{5.83} & 2.86 & 2.17 & 1.05 & 3.94 \\
Mean TTC [s] $\uparrow$       & \textbf{8.52} & 4.93 & 4.31 & 2.86 & 6.17 \\
PFR $\uparrow$                & \textbf{1.000} & 0.983 & 0.974 & 0.961 & 0.992 \\
\bottomrule
\end{tabular}
\vspace{1mm}

\begin{minipage}{\columnwidth}
\scriptsize
\emph{Notes:} Full, full method; No-GG, without game guidance; No-FSM, without
the finite-state machine; Fixed-C, fixed corridor. OSR loss is the
percentage-point difference from Full. PFR is the fraction of planning cycles with a feasible SOCP
solution. Bold denotes the best observed value.
\end{minipage}
\end{table}

Table~\ref{tab:ablation} shows that continuous corridor adaptation produced the
largest measured effect. Fixing the corridor reduced OSR by 18.6 percentage
points, increased contacts from 0 to 31, and gave the lowest PFR. Removing the
FSM reduced OSR by 12.4 points and
introduced 12 contacts, consistent with corridor revisions that are not held
long enough for the receding-horizon executor. Removing game guidance reduced
OSR by 9.3 points and increased TTO by 1.56~s. Thus the environment reward
alone did not reproduce the timing and clearance obtained with the
IBR-consistent prior and auxiliary target. Replacing TQC with SAC produced the
smallest OSR reduction, 6.2 points, but introduced five contacts. This outcome
is consistent with, but does not by itself prove, the proposed role of
truncated return quantiles in avoiding low-probability adverse corridors.

Two-sided count-level proportion tests with Holm adjustment across the four
ablations gave adjusted $p<10^{-6}$ for OSR in every case. As in the baseline
comparison, these values are descriptive because only aggregate counts are
included in the present statistical record. The ablation effect sizes,
collision counts, and PFR changes jointly support the conclusion that the
tactical corridor representation is the main source of the measured gain. The
ordering of the drops maps onto the proposed mechanism: adaptive corridor
shaping creates the opportunity, the FSM preserves it long enough to execute,
game guidance improves its timing and clearance, and TQC controls optimistic
distributional updates.
